% Part: propositional-logic
% Chapter: syntax-and-semantics
% Section: introduction

\documentclass[../../../include/open-logic-section]{subfiles}

\begin{document}

\olfileid[tr]{pl}{syn}{int}

\olsection{Giriş}

Önerme mantığı, !!{formula}s ile ilgilenir; bunlar !!{propositional variable}s
ve önerme bağlaçları $\lnot$, $\land$, $\lor$, $\lif$ ve $\liff$
kullanılarak kurulur. Sezgisel olarak !!a{propositional variable}~$p$, doğru
ya da yanlış olan bir tümceyi veya önermeyi temsil eder.
!!{propositional variable} terimi, !!a{formula} içinde geçen her değişken
için kullanılır. Bu değişkenlerin doğruluk değerleri belirlendiğinde,
onlardan önerme bağlaçlarıyla oluşturulan !!{formula}s için de doğruluk
değeri belirlenmiş olur. Önerme mantığının \emph{doğruluk
  fonksiyonlu} olduğunu söyleriz; çünkü semantiği doğruluk değerlerinin
fonksiyonlarıyla verilir. Özellikle önerme mantığında doğruluk ve yanlışlığın
daha ileri belirlenimlerini dikkate almayız: örneğin bir şeyin yalnızca
olumsal olarak değil zorunlu olarak doğru olup olmadığını, doğru olduğunun
bilinip bilinmediğini ya da şimdi doğru olmak yerine geçmişte doğru olmuş
veya gelecekte doğru olacak olup olmadığını ele almayız. Yalnızca doğru
($\True$) ve yanlış ($\False$) olmak üzere iki doğruluk değerini dikkate
alırız; dolayısıyla bir ifadenin ne doğru ne yanlış ya da yalnızca yarı doğru
olabilmesi olasılığını tartışma dışında bırakırız. Ayrıca yalnızca, bunlarla
kurulan !!a{formula} için doğruluk değerinin parçalarının doğruluk değerlerince
(örneğin anlamınca değil) bütünüyle belirlendiği bağlaçlara odaklanırız.
Özellikle, doğal dildeki koşulların doğruluk değerinin bu anlamda doğruluk
fonksiyonlu olup olmadığı tartışmalıdır. Maddi koşul~$\lif$ doğruluk
fonksiyonludur; başka mantıklar doğruluk fonksiyonlu olmayan koşulları ele
alır.

Doğruluk fonksiyonlu önerme mantığının kuramını ve metakuramını geliştirmek
için önce ifadelerinin sözdizimini ve semantiğini tanımlamamız gerekir.
!!{formula}s oluşturmanın bir yolunu betimleyeceğiz; bunu
!!{propositional variable}s ve bağlaçları kullanarak yapacağız. Başka tanımlar da mümkündür. Başka
sistemler farklı simgeler seçer, ilkel olarak farklı bağlaç kümelerini alır ve
parantezleri farklı kullanır (hatta sözde Leh gösteriminde olduğu gibi hiç
kullanmayabilir). Bununla birlikte bütün yaklaşımların ortak yanı, oluşum
kurallarının !!{formula}s kümesini \emph{tümevarımlı olarak} tanımlamasıdır.
Bu iş doğru yapılırsa her ifade, oluşum kurallarına göre özünde yalnızca tek
bir yoldan elde edilebilir. İfadelerin \emph{tek türlü okunabilir} olmasını
sağlayan tümevarımlı tanım, aynı yöntemi---tümevarımlı tanımı---kullanarak bu
ifadelere anlam verebilmemiz demektir.

İfadelere anlam vermek semantiğin alanıdır. Önerme mantığı semantiğindeki
merkezî kavram, !!a{valuation} altında sağlanmadır. !!^a{valuation}~$\pAssign{v}$,
her bir !!{propositional variable} için $\True$ ya da $\False$ doğruluk
değerlerinden birini belirler.
Her !!{valuation}, her !!{formula}~$!A$ için bir $\pValue{v}(!A)$ doğruluk
değeri belirler. !!^a{formula}, !!a{valuation}~$\pAssign{v}$ altında ancak ve
ancak $\pValue{v}(!A) = \True$ ise sağlanır---bunu $\pSat{v}{!A}$ biçiminde
yazarız. Bu bağıntı, mantıksal bağlaçların doğruluk fonksiyonlarından
yararlanarak, örneğin $!A \land !B$'nin sağlanmasını $!A$ ve~$!B$'nin
sağlanması (ya da sağlanmaması) cinsinden tanımlamak üzere, $!A$'nın yapısı
üzerinde tümevarımla da tanımlanabilir.

!!{formula}s için $\pSat{v}{!A}$ sağlama bağıntısını temel alarak totoloji,
mantıksal gerektirme ve
sağlanabilirlik gibi temel semantik kavramları tanımlayabiliriz.
!!^a{formula}, her !!{valuation} onu sağlıyorsa, yani her $\pAssign{v}$ için
$\pValue{v}(!A) = \True$ ise bir totolojidir, $\Entails !A$. Bir
!!{formula} kümesi~$\Gamma$, $!A$'yı mantıksal olarak gerektirir,
$\Gamma \Entails !A$, ancak ve ancak $\Gamma$'nın her üyesini sağlayan her
!!{valuation} $!A$'yı da sağlar. Bir
!!{formula}s kümesi, bir !!{valuation} söz konusu !!{formula}s kümesinin
tüm üyelerini aynı anda sağlıyorsa sağlanabilirdir. !!{formula}s tümevarımlı
olarak tanımlandığı ve sağlama da, !!{formula}s tek tek ele alındığında,
her birinin yapısı üzerinde tümevarımla
tanımlandığı için semantiğimizin özelliklerini kanıtlamak ve tanımlanan
semantik kavramları birbiriyle ilişkilendirmek üzere tümevarımı
kullanabiliriz.


\end{document}
